% =====================================================================
% Miary jakości podsumowań lingwistycznych T_1 - T_10 (T_11 - TODO)
%
% Źródło wzorów (rozdz. 8.2-8.4):
%   Niewiadomski A., "Methods for the Linguistic Summarization of Data:
%   Applications of Fuzzy Sets and Their Extensions",
%   Akademicka Oficyna Wydawnicza EXIT, Warszawa 2008,
%   ISBN 978-83-60434-40-6.
%   Seria: Badania Systemowe, t. 60, IBS PAN.
%
% Konwencja oznaczeń (zgodna z monografią Niewiadomskiego):
%   D            -- baza danych (zbiór m krotek)
%   m            -- liczba krotek (rekordów) w bazie D
%   n            -- liczba atrybutów branych pod uwagę w sumaryzatorze
%   y_i, d_i     -- i-ty obiekt / krotka
%   V_j(y_i)     -- wartość j-tego atrybutu obiektu y_i
%   X_j          -- uniwersum j-tego atrybutu (przestrzeń rozważań)
%   S            -- sumaryzator (ogólnie); S = S_1 t/s ... t/s S_n
%   S_j          -- j-ty zbiór rozmyty wchodzący w skład sumaryzatora
%                  na uniwersum X_j
%   W            -- kwalifikator (ogólnie); W = W_{g_1} t/s ... t/s W_{g_x}
%   W_{g_j}      -- j-ty zbiór rozmyty kwalifikatora na uniwersum X_{g_j}
%   x            -- liczba atrybutów wchodzących w skład kwalifikatora
%   Q            -- kwantyfikator lingwistyczny (zbiór rozmyty)
%   X_Q          -- uniwersum kwantyfikatora ([0,1] dla względnego,
%                  R^+ ∪ {0} dla bezwzględnego)
%   mu_A         -- funkcja przynależności zbioru rozmytego A
%   |A|          -- miara zbioru rozmytego A
%                  (sigma-count = suma mu_A dla zbiorów skończonych,
%                   całka funkcji mu_A dla nośników nieprzeliczalnych)
%   supp(A)      -- nośnik zbioru A: { x : mu_A(x) > 0 }
%   in(A)        -- stopień rozmycia (degree of fuzziness):
%                  in(A) = |supp(A)| / |X_A|
%   t, s         -- t-norma (AND) i t-conorma (OR)
% =====================================================================

\section*{Miary jakości podsumowań lingwistycznych}

Rozdział referuje miary jakości podsumowań lingwistycznych zdefiniowane
w monografii \cite{niewiadomski2008}, rozdz.~8.2--8.4. Pierwszych pięć
miar ($T_1$--$T_5$) odpowiada klasycznym ujęciom Yagera oraz Kacprzyka,
Yagera i Zadrożnego rozważanym w monografii; miary $T_6$--$T_{10}$ są
oryginalną propozycją A.~Niewiadomskiego, sformułowaną w sekcji 8.4
,,New quality measures of linguistic summaries''. Numeracja $T_i$
i kolejność prezentacji są zgodne z tym źródłem.

% =====================================================================
% T_1
% =====================================================================
\subsection*{$T_1$ -- Stopień prawdziwości \\
\textit{(degree of truth)}}

\paragraph{Co mierzy.}
Najbardziej znana miara jakości, wprowadzona przez Yagera. Określa,
\emph{w jakim stopniu prawdziwe} jest zdanie podsumowujące w sensie
rachunku zdań kwantyfikowanych lingwistycznie Zadeha; interpretowana
jako ,,poziom zaufania'' do podsumowania.

\paragraph{Wzór -- pierwsza forma ,,$Q\ P\ \text{are/have}\ S$''
(bez kwalifikatora).}
\begin{equation}
\label{eq:T1-form1}
    T_1\bigl(Q\ P\ \text{are/have}\ S\bigr)
    \;=\;
    \mu_Q\!\left( \frac{r}{M} \right),
    \qquad
    r = \sum_{i=1}^{m} \mu_S(d_i),
\end{equation}
gdzie $M = m$ dla kwantyfikatora \emph{względnego} i $M = 1$ dla
\emph{bezwzględnego}. Funkcję $\mu_S$ dla sumaryzatora złożonego
$S = S_1\ \text{AND/OR}\ \ldots\ \text{AND/OR}\ S_n$ wyznacza się
operatorami t-normy $t$ (AND) i t-conormy $s$ (OR) na cylindrycznych
rozszerzeniach do produktu uniwersów:
\begin{equation}
\label{eq:T1-mu-compound}
    \mu_S(d_i) \;=\; \mu_{S_1}(d_i)\ t/s\ \mu_{S_2}(d_i)\ t/s\ \ldots\ t/s\ \mu_{S_n}(d_i).
\end{equation}

\paragraph{Wzór -- druga forma ,,$Q\ P\ \text{being/having}\ W\ \text{are/have}\ S$''
(z kwalifikatorem).}
\begin{equation}
\label{eq:T1-form2}
    T_1\bigl(Q\ P\ \text{being/having}\ W\ \text{are}\ S\bigr)
    \;=\;
    \mu_Q\!\left(
        \frac{\sum_{i=1}^{m}\bigl(\mu_S(d_i) \wedge \mu_W(d_i)\bigr)}
             {\sum_{i=1}^{m}\mu_W(d_i)}
    \right).
\end{equation}
W mianowniku stoi $\Sigma\text{-count}(W)$ -- liczone są tylko krotki
o niezerowej przynależności do kwalifikatora. W tej formie
\textbf{możliwy jest tylko kwantyfikator względny}, ponieważ wartość
ułamka jest z natury proporcją.

\paragraph{Interpretacja.}
$T_1 \in [0, 1]$. Im wyższa wartość, tym podsumowanie jest bardziej
prawdziwe (informatywne). Wartość $T_1 = 0$ oznacza zdanie zupełnie
nietrafne, a $T_1 = 1$ -- w pełni potwierdzone przez zawartość bazy.

% =====================================================================
% T_2
% =====================================================================
\subsection*{$T_2$ -- Stopień nieprecyzyjności \\
\textit{(degree of imprecision)}}

\paragraph{Co mierzy.}
Ocenia, jak nieprecyzyjny jest \emph{sumaryzator}, na podstawie
stopnia rozmycia ($\mathrm{in}$) zbiorów rozmytych go tworzących.
Im szerszy nośnik zbiorów $S_j$, tym bardziej rozmyty (ogólny)
sumaryzator.

\paragraph{Wzór.}
\begin{equation}
\label{eq:T2}
    T_2 \;=\; 1 - \left( \prod_{j=1}^{n} \mathrm{in}(S_j) \right)^{1/n},
    \qquad
    \mathrm{in}(S_j) = \frac{|\mathrm{supp}(S_j)|}{|\mathcal{X}_j|}.
\end{equation}
Dla sumaryzatora prostego ($n=1$) wzór upraszcza się do
$T_2 = 1 - \mathrm{in}(S)$.

\paragraph{Interpretacja.}
$T_2 \in [0, 1]$. \textbf{Im wyższa wartość, tym precyzyjniejszy
(węższy) sumaryzator}, czyli lepsze podsumowanie. $T_2$ to dopełnienie
średniej geometrycznej stopni rozmycia -- gdy $\mathrm{in}(S_j)$
rośnie, $T_2$ maleje.

% =====================================================================
% T_3
% =====================================================================
\subsection*{$T_3$ -- Stopień pokrycia \\
\textit{(degree of covering)}}

\paragraph{Co mierzy.}
Mówi, jaka część krotek pasujących (choćby częściowo) do kwalifikatora
$W$ jest jednocześnie pokryta przez sumaryzator $S$. Pierwotnie
zdefiniowana dla podsumowań drugiej formy.

\paragraph{Wzór.}
\begin{equation}
\label{eq:T3}
    T_3 \;=\; \frac{\sum_{i=1}^{m} t_i}{\sum_{i=1}^{m} h_i},
\end{equation}
\begin{equation}
\label{eq:T3-ti}
    t_i = \begin{cases} 1, & \mu_S(d_i) > 0 \;\wedge\; \mu_W(d_i) > 0 \\ 0, & \text{w.p.p.}\end{cases},
    \qquad
    h_i = \begin{cases} 1, & \mu_W(d_i) > 0 \\ 0, & \text{w.p.p.}\end{cases}.
\end{equation}
Liczone są ostre wskaźniki należenia do nośnika, nie wartości $\mu$.

Dla pierwszej formy (bez kwalifikatora) traktuje się ją jako szczególny
przypadek, w którym $W = \mathcal{D}$ (,,być w bazie''), zatem
$h_i^{*} = 1$ dla wszystkich krotek, co daje:
\begin{equation}
\label{eq:T3-star}
    T_3^{*} \;=\; \frac{|\{ i : \mu_S(d_i) > 0\}|}{m} \;=\; \frac{|\mathrm{supp}(S \cap \mathcal{D})|}{m}.
\end{equation}

\paragraph{Interpretacja.}
$T_3 \in [0, 1]$. Im wyższa wartość, tym więcej rekordów wpadających
do kwalifikatora jest opisanych przez sumaryzator. Bardzo niskie
wartości oznaczają, że zdanie odnosi się do pomijalnie małej części
populacji (kwalifikator filtruje dużo, sumaryzator pokrywa mało).

% =====================================================================
% T_4
% =====================================================================
\subsection*{$T_4$ -- Stopień trafności \\
\textit{(degree of appropriateness)}}

\paragraph{Co mierzy.}
Sprawdza, czy podsumowanie jest \emph{nietrywialne} -- czy mówi coś,
czego nie dałoby się przewidzieć z samych częstości brzegowych
poszczególnych etykiet. Porównuje rzeczywiste pokrycie $T_3$
z iloczynem brzegowych częstości składowych sumaryzatora
(oczekiwanym pokryciem przy ,,niezależności'' cech).

\paragraph{Wzór.}
\begin{equation}
\label{eq:T4-rj}
    r_j \;=\; \frac{\sum_{i=1}^{m} g_{ij}}{m},
    \qquad
    g_{ij} = \begin{cases} 1, & \mu_{S_j}\bigl(V_j(y_i)\bigr) > 0 \\ 0, & \text{w.p.p.}\end{cases},
\end{equation}
\begin{equation}
\label{eq:T4}
    T_4 \;=\; \left| \prod_{j=1}^{n} r_j \;-\; T_3 \right|.
\end{equation}
$r_j$ -- udział rekordów spełniających pojedynczą $j$-tą etykietę
sumaryzatora; $\prod_j r_j$ -- ,,oczekiwane'' pokrycie przy założeniu
niezależności składowych.

\paragraph{Interpretacja.}
$T_4 \in [0, 1]$. \textbf{Im większa wartość, tym bardziej trafne
(odkrywcze) podsumowanie} -- pokrycie znacząco różni się od
oczekiwanego, co wskazuje na korelację cech ujętą przez konkretny
spójnik AND/OR. $T_4$ bliskie zera oznacza, że pokrycie jest dokładnie
takie, jak gdyby cechy były niezależne -- zdanie nie odkrywa nowej
informacji.

% =====================================================================
% T_5
% =====================================================================
\subsection*{$T_5$ -- Długość podsumowania \\
\textit{(length of summary)}}

\paragraph{Co mierzy.}
Karze podsumowania zbyt rozbudowane: im więcej zbiorów rozmytych
sklejonych spójnikami AND/OR tworzy sumaryzator $S$, tym mniej
czytelne (i mniej wartościowe) staje się zdanie.

\paragraph{Wzór.}
\begin{equation}
\label{eq:T5}
    T_5 \;=\; 2 \cdot (0{,}5)^{|S|},
\end{equation}
gdzie $|S|$ to \emph{liczba zbiorów rozmytych}, z których złożony jest
sumaryzator (uwaga: tu $|S|$ oznacza licznik składowych, a~nie miarę
zbioru rozmytego). Wartości:
$|S|=1 \Rightarrow T_5=1$,\;
$|S|=2 \Rightarrow T_5=0{,}5$,\;
$|S|=3 \Rightarrow T_5=0{,}25$, itd.

\paragraph{Interpretacja.}
$T_5 \in (0, 1]$. \textbf{Im wyższa wartość, tym krótsze (zwięźlejsze)
podsumowanie}, czyli lepsze. Nie istnieją wartości spoza tego
przedziału -- $T_5$ zawsze ma postać $2^{1-|S|}$.

% =====================================================================
% T_6
% =====================================================================
\subsection*{$T_6$ -- Stopień nieprecyzyjności kwantyfikatora \\
\textit{(degree of quantifier imprecision)}}

\paragraph{Co mierzy.}
Pierwsza z miar oryginalnie zaproponowanych przez Niewiadomskiego
\cite{niewiadomski2008}, sekcja~8.4.1. Analogiczna do $T_2$, ale
dotyczy kształtu zbioru rozmytego $Q$ reprezentującego \emph{kwantyfikator},
a nie sumaryzator. Im szerszy nośnik $Q$, tym mniej precyzyjny opis
ilościowy.

\paragraph{Wzór.}
\begin{equation}
\label{eq:T6}
    T_6 \;=\; 1 - \mathrm{in}(Q) \;=\; 1 - \frac{|\mathrm{supp}(Q)|}{|\mathcal{X}_Q|}.
\end{equation}
Mianownik sprowadza się do dwóch przypadków:
\begin{itemize}
    \item \textbf{kwantyfikator względny:} $|\mathcal{X}_Q| = 1$,
    bo $\mathcal{X}_Q = [0, 1]$;
    \item \textbf{kwantyfikator bezwzględny:} $|\mathcal{X}_Q| = m$,
    czyli liczba krotek w bazie.
\end{itemize}

\paragraph{Interpretacja.}
$T_6 \in [0, 1]$. \textbf{Im wyższa wartość, tym węższy nośnik
kwantyfikatora}, czyli tym bardziej precyzyjny opis ,,ile''. Dla
kwantyfikatorów ostrych (np.~\textsc{exactly~30} -- singleton
$\{1{,}0/30\}$, lub \textsc{for all} -- $\{1{,}0/1\}$) zachodzi
$T_6 = 0$, co intuicyjnie oznacza ich pełną precyzję na własnym
uniwersum.

% =====================================================================
% T_7
% =====================================================================
\subsection*{$T_7$ -- Stopień kardynalności kwantyfikatora \\
\textit{(degree of quantifier cardinality)}}

\paragraph{Co mierzy.}
Druga miara Niewiadomskiego \cite{niewiadomski2008}, sekcja~8.4.2.
Inspirowana obserwacją, że im większa kardynalność $|Q|$, tym mniej
precyzyjny kwantyfikator. \emph{Subtelnie różna od~$T_6$}: $T_6$ używa
miary nośnika $Q$, $T_7$ -- miary samego zbioru rozmytego, dzięki
czemu rozróżnia kwantyfikatory o tym samym nośniku, lecz różnym
kształcie (np.~trójkąt vs.~trapez na tym samym przedziale).

\paragraph{Wzór.}
\begin{equation}
\label{eq:T7}
    T_7 \;=\; 1 - \frac{|Q|}{|\mathcal{X}_Q|},
\end{equation}
gdzie $|Q| = \int \mu_Q(x)\,\mathrm{d}x$ (miara całkowa funkcji
przynależności) lub $\sum_x \mu_Q(x)$ (sigma-count) dla zbiorów
skończonych. Mianownik tak samo jak w $T_6$:
$|\mathcal{X}_Q| = 1$ (względny) lub $m$ (bezwzględny).

\paragraph{Interpretacja.}
$T_7 \in [0, 1]$. Im wyższa wartość, tym mniejsza miara $|Q|$, czyli
,,smuklejszy'' kwantyfikator -- bardziej precyzyjny.

% =====================================================================
% T_8
% =====================================================================
\subsection*{$T_8$ -- Stopień kardynalności sumaryzatora \\
\textit{(degree of summarizer cardinality)}}

\paragraph{Co mierzy.}
Trzecia miara Niewiadomskiego \cite{niewiadomski2008}, sekcja~8.4.3.
Ocenia łączną ,,smukłość'' sumaryzatora przez średnią geometryczną
znormalizowanych miar zbiorów rozmytych $S_1, \ldots, S_n$. Im mniejsze
miary $|S_j|$ względem swoich uniwersów, tym precyzyjniej etykiety
opisują wartości atrybutów.

\paragraph{Wzór.}
\begin{equation}
\label{eq:T8}
    T_8 \;=\; 1 - \left( \prod_{j=1}^{n} \frac{|S_j|}{|\mathcal{X}_j|} \right)^{1/n}.
\end{equation}

\paragraph{Interpretacja.}
$T_8 \in [0, 1]$. Im wyższa wartość, tym węższe (selektywniejsze)
poszczególne etykiety na swoich uniwersach. \textbf{Uwaga:}
normalizacja zachodzi przez \emph{uniwersum atrybutu}
$|\mathcal{X}_j|$, a~nie przez liczbę rekordów w bazie. $T_8$ ocenia
szerokość zbiorów na osi atrybutu, a~nie ich realną kardynalność
,,w bazie''.

% =====================================================================
% T_9
% =====================================================================
\subsection*{$T_9$ -- Stopień nieprecyzyjności kwalifikatora \\
\textit{(degree of qualifier imprecision)}}

\paragraph{Co mierzy.}
Czwarta miara Niewiadomskiego \cite{niewiadomski2008}, sekcja~8.4.4.
Określona dla podsumowań drugiej formy ,,$Q\ P\ \text{being}\ W\
\text{are}\ S$''. Analogon $T_2$ i~$T_6$: ocenia szerokość nośników
zbiorów rozmytych tworzących \emph{kwalifikator} $W$.

\paragraph{Wzór -- kwalifikator prosty.}
\begin{equation}
\label{eq:T9-single}
    T_9 \;=\; 1 - \mathrm{in}(W),
\end{equation}
gdzie $\mathrm{in}(W) = |\mathrm{supp}(W)|/|\mathcal{X}_W|$.

\paragraph{Wzór -- kwalifikator złożony z $W_{g_1}, \ldots, W_{g_x}$
na uniwersach $\mathcal{X}_{g_1}, \ldots, \mathcal{X}_{g_x}$.}
\begin{equation}
\label{eq:T9-multi}
    T_9 \;=\; 1 - \left( \prod_{j=1}^{x} \mathrm{in}(W_{g_j}) \right)^{1/x}.
\end{equation}

\paragraph{Interpretacja.}
$T_9 \in [0, 1]$. Im wyższa wartość, tym węższy (precyzyjniej
zdefiniowany) kwalifikator -- zdanie filtruje szczegółowo określoną
podgrupę obiektów.

% =====================================================================
% T_10
% =====================================================================
\subsection*{$T_{10}$ -- Stopień kardynalności kwalifikatora \\
\textit{(degree of qualifier cardinality)}}

\paragraph{Co mierzy.}
Piąta miara Niewiadomskiego \cite{niewiadomski2008}, sekcja~8.4.5.
Bazuje na kardynalności (mierze) zbiorów $W_{g_j}$. Względem $T_9$
relacja taka sama jak $T_7$ względem $T_6$: rozróżnia kwalifikatory
o tym samym nośniku, lecz różnym kształcie.

\paragraph{Wzór -- kwalifikator prosty.}
\begin{equation}
\label{eq:T10-single}
    T_{10} \;=\; 1 - \frac{|W|}{|\mathcal{X}_g|}.
\end{equation}

\paragraph{Wzór -- kwalifikator złożony.}
\begin{equation}
\label{eq:T10-multi}
    T_{10} \;=\; 1 - \left( \prod_{j=1}^{x} \frac{|W_{g_j}|}{|\mathcal{X}_{g_j}|} \right)^{1/x}.
\end{equation}

\paragraph{Interpretacja.}
$T_{10} \in [0, 1]$. Im wyższa wartość, tym mniejsza miara
$|W_{g_j}|$ względem $|\mathcal{X}_{g_j}|$ -- ,,smuklejszy''
kwalifikator. \textbf{Uwaga:} tak jak $T_8$, miara normalizowana jest
przez uniwersum atrybutu, a~nie przez liczbę rekordów w bazie. Nie
mówi więc, ilu rekordów dotyczy kwalifikator, tylko jak wąsko jest
on określony na osi swojego atrybutu.

% =====================================================================
% Goodness of summary (8.74)
% =====================================================================
\subsection*{Łączna jakość podsumowania \\
\textit{(goodness of the summary)}}

Łączną miarę jakości definiuje się jako średnią ważoną składowych
miar~$T_i$ \cite[wzór~(8.74)]{niewiadomski2008}:
\begin{equation}
\label{eq:goodness}
    T \;=\; T(T_1, \ldots, T_{10};\ w_1, \ldots, w_{10})
    \;=\; \sum_{i=1}^{10} w_i \cdot T_i,
    \qquad
    \sum_{i=1}^{10} w_i = 1.
\end{equation}
Optymalne podsumowanie wybiera się jako:
\begin{equation}
\label{eq:goodness-argmax}
    S^{*} \;=\; \arg\max_{S \in \mathbb{S}}\; \sum_{i=1}^{10} w_i \cdot T_i.
\end{equation}
Autor zaleca (gdy nic innego nie wpływa na wybór wag) ustawienie
$w_1 = w_2 = w_4 = w_5 = w_6 = w_7 = w_8 (= w_3)$ dla podsumowań
pierwszej formy oraz $w_1 = \ldots = w_{10}$ dla podsumowań drugiej
formy (jednakowe wagi dla wszystkich miar).

% =====================================================================
% T_11 - placeholder
% =====================================================================
\subsection*{$T_{11}$ -- Długość kwalifikatora \\
\textit{(length of qualifier)}}

\textbf{TODO.} Miara nie pojawia się w monografii Niewiadomskiego
\cite{niewiadomski2008} -- rozdz.~8.4 i~zbiorczy wzór~(8.74) obejmują
wyłącznie $T_1$--$T_{10}$. $T_{11}$ jest najpewniej rozszerzeniem
narzuconym przez prowadzącego (analogon $T_5$ dla kwalifikatora).
Spodziewana postać przez analogię:
\[
    T_{11} \;=\; 2 \cdot (0{,}5)^{|W|},
\]
gdzie $|W|$ to liczba zbiorów rozmytych tworzących kwalifikator.
Po znalezieniu wzoru w treści projektu lub materiałach prowadzącego
sekcję uzupełnić.

% =====================================================================
% Bibliografia (do scalenia z głównym sprawozdaniem)
% =====================================================================
% \bibitem{niewiadomski2008}
% Niewiadomski, A.: \textit{Methods for the Linguistic Summarization
% of Data: Applications of Fuzzy Sets and Their Extensions}.
% Akademicka Oficyna Wydawnicza EXIT, Warszawa 2008,
% ISBN~978-83-60434-40-6. Seria: Badania Systemowe, t.~60,
% Instytut Badań Systemowych PAN.
